\chapter{Discussion}

In this thesis I first reproduced two behavioral models (BFE- and GFE agent) from ... that are capable of perception, planning, and action. Both are based on free energy functionals from which message passing schemes are derived. Message passing enables them with inference, which in turn gives rise to their full behavioral repetoire. The GFE agent is based on the BFE agent but additionally entails a mechanism that reduces the uncertainty of its hidden state (part of its generative model) by seeking information in the environment via actions, i.e., it is curious.
In a second step, I expanded the free energy functional of the BFE agent to introduce a hidden context state in the generative model, which gave rise to the Context-BFE agent. The corresponding message passing algorithm was derived from this expanded free energy.
In a third step, I took the integration of the information seeking mechanism in the GFE agent as an example to introduce this mechanism to the Context-BFE agent for the context state, i.e., the mechanism leads to information seeking that reduces uncertainty of the hidden context state. This integration finally led to the Context-GFE agent. From its customized free energy functional I derived the corresponding message passing algorithm.
I tested all four agents in the simulated T-Maze environment which is frequently used to reveal information seeking behavior. There I showed that the newly developed Context-GFE agent shows similar behavior to the GFE agent, which strongly suggests that the integration of the information seeking mechanism in the Context-GFE agent was successful, i.e., it is curious like the GFE agent and is driven to reduce uncertainty over its hidden (context) state. As a baseline, I also gave the results for both non-curious agents (BFE- and Context-BFE), which both show similar (non-information seeking) behavior.
Beyond that, I generally addressed the kind of behavior that results from the information seeking mechanism in both curious agents (not uniquely tied to the newly developed Context-GFE agent).

\section{Interpretation}



-alternative meaning of success
-attention parameter
-beta sweep
-alpha sweep for comparability with de Vries papers

\section{Limitations}

\subsubsection{Compromises of the Mean-Field Approximation in the Context-GFE agent}

While the results of both curious agents are qualitatively similar, they are not equal, meaning that although the Context-GFE agent also performs information seeking like the GFE agent, there are subtle differences that are addressed in the following. The Context-GFE agent shows a more abrupt transition of probabilities from smaller \(c\) values and \(\alpha\) values near \(0.5\) to larger \(c\) values and more extreme \(\alpha\) values. This phenomenon is most likely a consequence of a mean-field approximation that was necessary to develop the Context-GFE agent. Specifically, while the GFE agent models the T-Maze type as part of its hidden state, the Context-GFE agent instead models it via its context state. The differences between the two approaches are as follows. The hidden state is modeled individually for each timestep in the trial (simulation) while there is only a single context state that persists for the whole trial (see ...). Furthermore, all three observations are related to this single context state in the agent's generative model, i.e, in the factor graph representation there are edges between all observation states and the context state. This leads to cycles in the factor graph which can lead to inference algorithms that do no converge to a stable solution in the case when the statistical dependencies between variables at each node are respected by the local belief distributions, i.e, when we assume the Bethe free energy at each local node. As a remedy, it is sufficient to break these local statistical dependencies at single nodes to remove the cycles. This was done between the observation states and the context state and leads to mean-field approximations at these local nodes. As a result, the inference process leads to stable solutions but has the problem of disrespecting uncertainties in the generative model related to context- and observation states. As a consequence, inference can yield results that are overconfident and thus are not as nuanced as without using the mean-field approximation. This circumstance becomes evident in the results of both curious agents where the GFE agent does not rely on the mean-field approximation between its observation states and hidden states (it does not have a context state and instead models the T-Maze type via the hidden states) and thus achieves a nuanced behavioral shift with varying \(\alpha\)- and \(c\) values. In contrast, the Context-GFE agent relies on the mean-field approximation at the prescribed locations in its generative model and becomes overly confident more quickly with varying parameter combinations (\(\alpha\) and \(c\)) which leads to the visible abrupt change of probabilities for larger \(c\) values and extreme \(\alpha\) values.

\subsubsection{Alternative Meaning of Success}

This situation should not be mistaken to be a weakness of curious agents, instead it shows their inclination of not only preferring predicted observations (like non-curious agents do) but also value hidden environmental causes that are less ambiguous, i.e, they seek to reduce uncertainty about their model of these causes, which are the hidden states. This mechanism may allow them to reduce their variational free energy more reliably in the long term and thus better secure their existence in evolutionary terms (according to the Free Energy Principle). In this way, the term "Long-Term Success" might not only be understood as the rates of receiving the reward but additionally refers to the agent's ability to minimize its variational free energy which necessitates the curiousity mechanism.


\section{Outlook}